% =====================================================================
\section{Priming: the only route that reached the target}\label{sec:prime}
Priming is memory of a past exposure: a cell that has already met an oxidant responds to the next exposure more strongly and more quickly. It is inserted into the model in one line: the same signal that drives induction accumulates a trace $p$, and the induction strength is multiplied by $(1+A\,p)$. The ceiling $A$ is taken from measurement --- threefold.

\facts
\begin{itemize}
\item On a repeated exposure 24 hours later the response of the regulator is lower while the accumulation of the enzyme itself is \emph{higher} --- up to threefold --- and the time to half-maximum falls from about 8 to about 4 hours. The same laboratory and the same system that gave us the raw kinetic data.\tS
\item The effect depends on the master regulator of the oxidative response and does not depend on its auxiliary partner: that is, it is an amplification of the signal and not a side effect.\tS
\item Independently, in another tissue: accumulation of a protective enzyme under repeated exposures was followed over four days and was gene-selective.\tP
\end{itemize}

\textbf{What came out.} For the first time in the project the episode duration entered the clinical target: median 1.46 days against an upper bound of 1.665 (a one-sided estimate from order statistics), fraction of short episodes 0.363, paired test against the control without priming $p=5.2\cdot10^{-6}$. All four parts of the rule written down before the computation are satisfied --- and satisfied again on a sample of 2000 sets, six times larger.

It matters what this is \emph{not}. Priming was not fitted: its amplitude is taken as the ceiling from measurement, its time constant was scanned, and the acceleration of the first stage --- the second half of the measured effect --- was checked separately and gave almost nothing, exactly as predicted before the computation. One route remained, and it passed.
